\section*{Chapitre \ref{ch:b2:atomic-orbitals} --- Orbitales atomiques~: formes, énergies et tailles}

\begin{solution}{exo:b2:atomic-orbitals:1}
3p~: $n = 3$, $l = 1$~; un \omterm{def:b2:atomic-orbitals:node}{nœud radial} ($n - l - 1$) et un \omterm{def:b2:atomic-orbitals:node}{nœud angulaire}. 4d~: $n = 4$,
$l = 2$~; un \omterm{def:b2:atomic-orbitals:node}{nœud radial} et deux \omterm{def:b2:atomic-orbitals:node}{nœuds angulaires}.
\end{solution}

\begin{solution}{exo:b2:atomic-orbitals:2}
$P_{2p} \propto r^4\mathrm e^{-r}$ ($r$ en $a_0$)~: $d\ln P/dr = 4/r - 1 = 0$ pour $r = 4$~:
$4a_0 = \qty{212}{pm}$.
\end{solution}

\begin{solution}{exo:b2:atomic-orbitals:3}
$1 - \mathrm e^{-2}(1 + 2 + 2) = 1 - 5\mathrm e^{-2} = 0.323$.
\end{solution}

\begin{solution}{exo:b2:atomic-orbitals:4}
Quatre lobes pointant entre les axes $x$ et $y$, dans le plan $xy$, de signes
alternés (positifs là où $xy > 0$)~; les plans nodaux sont $xz$ et $yz$. \omterm{def:b2:atomic-orbitals:node}{Nœuds radiaux}~:
$n - l - 1 = 0$.
\end{solution}

\begin{solution}{exo:b2:atomic-orbitals:5}
Oxygène, $(1\mathrm s)^2(2\mathrm s, 2\mathrm p)^6$~: $\sigma = 5 \times 0.35 + 2 \times 0.85
= 3.45$, $Z^* = 4.55$~; rayon $4/4.55 = 0.88a_0$, \qty{47}{pm}.
\end{solution}

\begin{solution}{exo:b2:atomic-orbitals:6}
L'ionisation arrache les électrons les moins liés. Les électrons 4s
(\qty{-14}{eV}) sont bien moins liés que les électrons 3d (\qty{-59}{eV})~: les deux électrons 4s
partent d'abord, puis un électron 3d, ce qui donne $[\ce{Ar}]\,3\mathrm d^5$.
\end{solution}

\begin{solution}{exo:b2:atomic-orbitals:7}
Lithium, 2s~: $Z^* = 3 - 2 \times 0.85 = 1.30$, $I = 13.6 \times 1.30^2/4 =
\qty{5.75}{eV}$ (mesurée 5.39, trop haute de 7~\%). Sodium, 3s~: $Z^* = 11 - 8 \times 0.85 -
2 \times 1.00 = 2.20$, $I = 13.6 \times 2.20^2/9 = \qty{7.3}{eV}$ (mesurée 5.14)~: les
règles, modèle grossier, surestiment la liaison de l'électron externe, de plus en
plus en descendant le groupe.
\end{solution}

\begin{solution}{exo:b2:atomic-orbitals:8}
Sodium~: $Z^* = 2.20$, rayon $9/2.20 = 4.1a_0$ (\qty{216}{pm}). Chlore, 3p~:
$\sigma = 6 \times 0.35 + 8 \times 0.85 + 2 \times 1.00 = 10.90$, $Z^* = 6.10$, rayon
$9/6.10 = 1.5a_0$ (\qty{78}{pm}). Même couche, mais l'électron externe du chlore
perçoit une charge presque trois fois plus grande, et il est attiré vers le noyau.
\end{solution}

\begin{solution}{exo:b2:atomic-orbitals:9}
Toutes deux ont pour \omterm{def:b2:atomic-orbitals:radial-angular}{partie angulaire} $1/\sqrt{4\pi}$~: l'intégrale se réduit donc à
\[
  \int_0^\infty R_{1s}R_{2s}r^2\,dr \propto \int_0^\infty (2 - r)r^2\mathrm e^{-3r/2}\,dr =
  2 \times \frac{2!}{(3/2)^3} - \frac{3!}{(3/2)^4} = \frac{32}{27} - \frac{32}{27} = 0 .
\]
\end{solution}

\begin{solution}{exo:b2:atomic-orbitals:10}
$\int_0^\infty N^2r^2\mathrm e^{-r}r^2\,dr = N^2 \times 4! = 1$~: $N = 1/\sqrt{24} =
1/(2\sqrt6)$, comme dans le tableau.
\end{solution}

\begin{solution}{exo:b2:atomic-orbitals:11}
$\langle r\rangle = 1.5a_0$, \omterm{def:b2:atomic-orbitals:radial-distribution}{rayon le plus probable} $a_0$. $P(r) = 4r^2\mathrm e^{-2r}$ monte
brutalement à partir de zéro et décroît lentement~: la longue queue aux grands $r$ pèse davantage dans la
moyenne que dans la position du maximum.
\end{solution}

\begin{solution}{exo:b2:atomic-orbitals:12}
$Y_{p_z} \propto z/r = \cos\theta$ s'annule pour $\theta = \pi/2$~: le plan $xy$.
$\psi$ a le signe de $z$, positive au-dessus, négative au-dessous. Sur le plan, la
\omterm{def:b2:atomic-orbitals:wavefunction}{densité de probabilité} est nulle.
\end{solution}

\begin{solution}{pb:b2:atomic-orbitals:1}
\textbf{1.} $\int|\psi|^2\,dV = \frac{4\pi}{\pi a_0^3}\int_0^\infty r^2\mathrm e^{-2r/a_0}\,dr =
\frac{4}{a_0^3} \times \frac{a_0^3}{4} = 1$.
\textbf{2.} $P(r) = 4\pi r^2|\psi|^2 = (4/a_0^3)\,r^2\mathrm e^{-2r/a_0}$.
\textbf{3.} $dP/dr = 0$~: $2/r - 2/a_0 = 0$, $r = a_0$.
\textbf{4.} $\langle r\rangle = (4/a_0^3)\int_0^\infty r^3\mathrm e^{-2r/a_0}\,dr = (4/a_0^3)
\times 6(a_0/2)^4 = 1.5a_0$.
\textbf{5.} La distribution est dissymétrique, avec une longue queue aux grands $r$.
\textbf{6.} $|\psi|^2$ est une densité par unité de volume, maximale au noyau~; $P(r)$
tient compte du volume de la coquille, $4\pi r^2\,dr$, qui s'annule en $r = 0$.
\textbf{7.} \qty{52.9}{pm}.
\textbf{8.} Avec $x = r/a_0$~: $\int_{x_0}^\infty 4x^2\mathrm e^{-2x}\,dx$~; par deux intégrations par parties,
$= \mathrm e^{-2x_0}(2x_0^2 + 2x_0 + 1)$.
\textbf{9.} $5\mathrm e^{-2} = 0.677$.
\textbf{10.} $13\mathrm e^{-4} = 0.238$.
\textbf{11.} $\mathrm e^{-2x}(1 + 2x + 2x^2) = 0.10$ pour $x \approx 2.66$~: \qty{141}{pm}.
\textbf{12.} La densité ne s'annule à aucune distance~: pas de bord. Mais 90~\% de la
probabilité se trouvent à moins de \qty{141}{pm}~: c'est une taille pratique.
\textbf{13.} 2s~: un \omterm{def:b2:atomic-orbitals:node}{nœud radial} (sphère de rayon $2a_0$), aucun \omterm{def:b2:atomic-orbitals:node}{nœud angulaire}. 2p~: un
\omterm{def:b2:atomic-orbitals:node}{nœud angulaire} (un plan), aucun \omterm{def:b2:atomic-orbitals:node}{nœud radial}.
\textbf{14.} En $r = 2a_0$.
\textbf{15.} $4a_0$ (exercice 2).
\textbf{16.} $d\ln P/dr = 2/r - 2/(2 - r) - 1 = 0$ donne $r^2 - 6r + 4 = 0$, $r = 3 \pm \sqrt5$~:
$0.76a_0$ (petit maximum interne) et $5.24a_0$ (maximum principal).
\textbf{17.} 2s, grâce à son maximum interne. Dans un atome polyélectronique, un électron 2s
pénètre la couche 1s, est moins écranté qu'un électron 2p et se situe plus bas en
énergie.
\textbf{18.} Tous les rayons sont divisés par $Z = 6$~: 2p à $0.67a_0$ (\qty{35}{pm})~; maximums de 2s
à $0.13a_0$ et $0.87a_0$.
\textbf{19.} C 3.25, N 3.90, O 4.55.
\textbf{20.} $\varepsilon = -13.6\,Z^{*2}/4$~: C \qty{-35.9}{eV}, N \qty{-51.7}{eV}, O
\qty{-70.4}{eV}.
\textbf{21.} Comme dans le chapitre~: \qty{11.5}{eV} (mesurée 11.26).
\textbf{22.} N~: atome $5 \times (-13.6 \times 3.90^2/4) = \qty{-258.6}{eV}$~; \ce{N+}, $Z^* =
4.25$, $4 \times (-13.6 \times 4.25^2/4) = \qty{-245.7}{eV}$~; $I = \qty{12.9}{eV}$.
\textbf{23.} O~: atome $6 \times (-13.6 \times 4.55^2/4) = \qty{-422.3}{eV}$~; \ce{O+}, $Z^* =
4.90$, $5 \times (-13.6 \times 4.90^2/4) = \qty{-408.2}{eV}$~; $I = \qty{14.2}{eV}$. Modèle
11.5, 12.9, 14.2~; mesures 11.26, 14.53, 13.62.
\textbf{24.} L'énergie mesurée de l'oxygène est inférieure à celle de l'azote~: dans
l'oxygène, le quatrième électron 2p doit partager une orbitale avec un autre, et leur
répulsion le rend plus facile à arracher, tandis que les trois électrons p célibataires de l'azote
forment une sous-couche à demi remplie stable. Les \omterm{def:b2:atomic-orbitals:slater}{règles de Slater} traitent tous les électrons d'un groupe
de la même façon et ignorent tout de l'appariement.
\textbf{25.} L'électron 1s de l'hydrogène se trouve au-delà de $2a_0$ avec la probabilité
$\boldsymbol{13\mathrm e^{-4} \approx 0.238}$.
\end{solution}
